\chapter{Evaluation}
\label{ch:evaluation}

This chapter describes the experimental evaluation of the MRC system.
First, the simulation environments and test infrastructure are presented
(\autoref{sec:simulation_environment}).
The evaluation methodology, metrics, and automated test protocol are
then defined (\autoref{sec:evaluation_methodology}).
Subsequently, the experimental results are analyzed per fleet size:
3~robots (\autoref{sec:results_3robots}), 5~robots
(\autoref{sec:results_5robots}), 10~robots
(\autoref{sec:results_10robots}), and a realistic two-robot proximity
scenario (\autoref{sec:proximity_test}).
Finally, the cross-experiment impact of ICP refinement
(\autoref{sec:icp_impact}) and pose graph optimization
(\autoref{sec:pose_graph_impact}) is analyzed, followed by a summary
of key findings (\autoref{sec:evaluation_summary}).

%% ---------------------------------------------------------------------------
\section{Simulation Environment}
\label{sec:simulation_environment}

The evaluation was conducted in a physics-based simulation using Gazebo
(Harmonic) with the ROS\,2 Jazzy middleware.
A dedicated ROS\,2 package (\texttt{mrc\_test}) was developed to
provide reproducible, automated benchmarking across a wide range of
configurations and fleet sizes.
Two complementary simulation environments were used: a detailed
indoor room for functional validation and a larger, reduced-physics
environment for scalability testing.

\subsection{Functional Test World}
\label{sec:test_world}

The primary test environment consists of an $8 \times 8$\,m indoor
room with \SI{2.5}{\metre} walls.
To provide geometric features for LiDAR-based registration, the
room is furnished with several static objects of varying shapes
and sizes: a table ($1.2 \times 0.8$\,m), a chair, a bookshelf
($1.5 \times 0.4 \times 1.8$\,m), a cylindrical trash can
($r = 0.2$\,m), a crate ($0.6^3$\,m), a sphere ($r = 0.15$\,m),
a plant pot, and a square pillar ($0.3 \times 0.3 \times 2.5$\,m).
These objects are distributed asymmetrically to avoid ambiguities
in the registration process.

Three heterogeneous robots are deployed in an equilateral triangle
formation with a side length of \SI{1.5}{\metre}:
\begin{enumerate}
    \item \textbf{EC Swift~1} (master/reference robot):
          Positioned at $(-0.75, -0.433, 0.15)$\,m with yaw $= 0\degree$.
          Equipped with a single front-facing Livox Mid-360 LiDAR.
    \item \textbf{EC Swift~2}:
          Positioned at $(0.75, -0.433, 0.15)$\,m with yaw $= 120\degree$.
          Equipped with front and back Livox Mid-360 LiDARs, merged
          into a single point cloud.
    \item \textbf{Athena}:
          Positioned at $(0.0, 0.866, 0.15)$\,m with yaw $= -90\degree$.
          Equipped with front and back Livox Mid-360 LiDARs, merged
          into a single point cloud.
\end{enumerate}

\subsection{Scalability Test World}
\label{sec:scale_world}

A second simulation world was designed specifically for evaluating
MRC performance as the fleet size increases.
The scalability world is a $20 \times 20$\,m room with walls and
distributed furniture, using a reduced-physics configuration: only
the ODE physics engine (for stable placement), the sensor rendering
system, and the RGL LiDAR manager are loaded.
IMU and contact simulation are disabled to reduce the computational
overhead per robot and enable simulations with up to 20 simultaneous
robots.

All robots in the scalability tests are homogeneous EC~Swift platforms,
each equipped with a single front-facing Livox Mid-360 LiDAR.
Robot~0 always serves as the master.

\subsection{Predefined Position Configurations}
\label{sec:position_configs}

To ensure systematic and reproducible spatial evaluation, a
comprehensive YAML configuration file defines robot positions
for fleet sizes of $N \in \{2, 3, 5, 7, 10, 12, 15, 17, 20\}$
robots.
For each fleet size, five distinct position setups are provided,
yielding a total of 45 predefined spatial arrangements.
Each position entry specifies $(x, y, z, \psi)$ in meters and
radians.

The five setups per fleet size follow a systematic pattern
designed to exercise different aspects of registration performance:

\begin{description}
    \item[Setup~1: Tight regular polygon.]
        Robots are arranged as a regular $N$-gon with a small
        radius (e.g., \SI{2}{\metre} for 3~robots,
        \SI{5}{\metre} for 20~robots), maximizing inter-robot
        LiDAR overlap.
    \item[Setup~2: Wide regular polygon.]
        The same $N$-gon with a larger radius (e.g., \SI{4}{\metre}
        for 3~robots, \SI{8}{\metre} for 20~robots), testing
        performance under reduced overlap conditions.
    \item[Setup~3: Line/grid formation.]
        Robots are placed in a structured grid or parallel-line
        arrangement (e.g., two rows of five for 10~robots,
        a $5 \times 4$ grid for 20~robots).
    \item[Setup~4: Structured composite.]
        More complex regular patterns such as concentric polygons
        or a center robot with surrounding rings.
    \item[Setup~5: Scattered asymmetric.]
        Robots are placed at irregular, non-symmetric positions
        with varying orientations to test robustness under
        realistic conditions.
\end{description}

Each robot's heading ($\psi$) is varied across setups to create
diverse fields of view.
All positions remain within \SI{8.5}{\metre} of the origin to
fit within the $20 \times 20$\,m room.

\subsection{LiDAR Simulation}

The Livox Mid-360 LiDAR is simulated using the RGL (Robotec GPU
Lidar) Gazebo plugin, which reproduces the non-repetitive scanning
pattern of solid-state LiDARs.
This is a relevant design choice, as the non-repetitive scan pattern
requires temporal accumulation for sufficient spatial coverage, which
is handled by the slave nodes' time-window accumulation mode with
$\Delta t_\mathrm{acc} = \SI{2.0}{\second}$.

%% ---------------------------------------------------------------------------
\section{Evaluation Methodology}
\label{sec:evaluation_methodology}

\subsection{Evaluator Node}
\label{sec:evaluator_node}

An automated evaluator node (\texttt{mrc\_evaluator}, approximately
1\,600 lines of C++) was implemented to enable reproducible
benchmarking.
The evaluator operates within the same TF tree as the master node and
compares the calibrated transforms against the ground truth after each
calibration run.
Key features of the evaluator include:

\begin{itemize}
    \item \textbf{Dynamic robot discovery:}
          The evaluator subscribes to the same
          \texttt{/robot\_announcement} topic as the master node.
          When a new \texttt{RobotAnnouncement} message is received,
          the evaluator automatically registers the robot, loads its
          ground truth from launch-file parameters, and creates per-robot
          error publishers.

    \item \textbf{Ground truth from spawn positions:}
          Ground truth poses are derived from the known Gazebo spawn
          positions, transformed into the master robot's
          \texttt{base\_link} frame:
          \begin{equation}
              \mathbf{p}_\mathrm{gt}^{(i)} = \mathbf{R}(-\psi_0)
              \cdot (\mathbf{p}_i - \mathbf{p}_0), \qquad
              \psi_\mathrm{gt}^{(i)} = \psi_i - \psi_0
          \end{equation}
          where $(\mathbf{p}_0, \psi_0)$ is the master's world pose
          and $(\mathbf{p}_i, \psi_i)$ is the $i$-th slave's world pose.

    \item \textbf{Voxel-based overlap computation:}
          The overlap ratio for slave $i$ is:
          \begin{equation}
              o^{(i)} = \frac{|V_\mathrm{master} \cap V_i|}{|V_i|}
          \end{equation}
          where $V_\mathrm{master}$ and $V_i$ are the occupied voxel
          sets of the master and slave point clouds at resolution
          $v = \SI{0.25}{\metre}$.

    \item \textbf{Runtime parameter reconfiguration:}
          The evaluator modifies MRC master parameters (algorithm,
          pose graph mode, topology, ICP refinement) at runtime via
          the ROS\,2 parameter service, enabling all 12 configurations
          to be tested within a single simulation session.
\end{itemize}

\subsection{Metrics}
\label{sec:metrics}

For each calibrated robot $i$, the following error metrics are computed
against the ground truth:

\paragraph{Translation error.}
The Euclidean distance between the estimated and ground truth positions:
\begin{equation}
    e_t^{(i)} = \lVert \mathbf{p}_\mathrm{est}^{(i)}
                      - \mathbf{p}_\mathrm{gt}^{(i)} \rVert_2
\end{equation}

\paragraph{Rotation error.}
The absolute yaw angle difference, normalized to $[-\pi, \pi]$:
\begin{equation}
    e_r^{(i)} = \left| \psi_\mathrm{est}^{(i)}
                     - \psi_\mathrm{gt}^{(i)} \right|
\end{equation}

These per-robot errors are aggregated into the following summary
statistics over $n$ valid robots:
\begin{itemize}
    \item \textbf{Mean translation error:}
          $\bar{e}_t = \frac{1}{n} \sum_{i=1}^{n} e_t^{(i)}$
    \item \textbf{RMS translation error:}
          $e_t^\mathrm{RMS} = \sqrt{\frac{1}{n} \sum_{i=1}^{n}
          (e_t^{(i)})^2}$
    \item \textbf{Maximum translation error:}
          $e_t^\mathrm{max} = \max_i \, e_t^{(i)}$
    \item \textbf{Mean rotation error:}
          $\bar{e}_r = \frac{1}{n} \sum_{i=1}^{n} e_r^{(i)}$
    \item \textbf{Mean overlap ratio:}
          $\bar{o} = \frac{1}{n} \sum_{i=1}^{n} o^{(i)}$
\end{itemize}

\subsection{Test Protocol and Automation}
\label{sec:test_protocol}

Each configuration was evaluated using $k = 3$ independent calibration
trials.
The reported metrics are averaged across trials to reduce the influence
of stochastic effects (e.g., variations in accumulated point clouds
due to the non-repetitive LiDAR pattern).

The configurations tested span the full combinatorial space of the
MRC system parameters:
\begin{itemize}
    \item \textbf{Registration backend:} KISS-Matcher, FRICP
    \item \textbf{Pose graph optimization:} disabled, star topology,
          mesh topology
    \item \textbf{ICP refinement:} disabled, enabled
\end{itemize}
This yields $2 \times 3 \times 2 = 12$ configurations, each
evaluated with $k$ trials, for a total of $12k$ calibration runs
per test setup.

\paragraph{Automated test suite.}
A Bash orchestration script (\texttt{run\_all\_mrc\_tests.bash})
automates the complete evaluation across multiple robot counts and
position setups.
The default test matrix comprises $3$ robot counts $\times$
$5$ position setups $= 15$ tests, each evaluating 12~configurations
with $k = 3$~trials, for a total of $15 \times 12 \times 3 = 540$
calibration runs.
All results are written to a timestamped output directory.

%% ---------------------------------------------------------------------------
\section{Results: 3-Robot Experiments}
\label{sec:results_3robots}

Three heterogeneous robots (2~EC~Swift + 1~Athena) were deployed in
the functional test world ($8 \times 8$\,m room) across 5 position
setups.
All 5 tests completed successfully within the \SI{1000}{\second}
timeout, with test durations between 477 and \SI{498}{\second}.
The mean point cloud overlap across setups was 58.6\%
(range: 51.4--64.2\%).

\autoref{tab:results_3r} presents the averaged results across all
5 setups and 3 trials per configuration.

\begin{table}[htbp]
    \centering
    \caption{3-robot calibration results (averaged over 5~position
             setups, 3~trials per configuration, overlap: 58.6\%).}
    \label{tab:results_3r}
    \small
    \begin{tabular}{l r r r r}
        \toprule
        \textbf{Configuration}
            & $\bar{e}_t$ [m]
            & $e_t^\mathrm{RMS}$ [m]
            & $\bar{e}_r$ [deg]
            & $t$ [s] \\
        \midrule
        \multicolumn{5}{l}{\textit{KISS-Matcher}} \\
        \quad no PG             & 2.97  & 3.44  & 49.5  & 0.1 \\
        \quad star PG           & 2.74  & 3.19  & 48.6  & 0.1 \\
        \quad mesh PG           & 2.55  & 2.85  & 44.6  & 0.1 \\
        \quad no PG + ICP       & \textbf{2.04}  & \textbf{2.37}  & 48.8  & 1.2 \\
        \quad star PG + ICP     & 2.68  & 3.27  & 70.1  & 1.2 \\
        \quad mesh PG + ICP     & 2.81  & 3.31  & 73.6  & 1.6 \\
        \midrule
        \multicolumn{5}{l}{\textit{FRICP}} \\
        \quad no PG             & 4.06  & 4.43  & 120.3 & 1.9 \\
        \quad star PG           & 4.42  & 4.75  & 133.1 & 1.9 \\
        \quad mesh PG           & 5.11  & 5.32  & 138.7 & 2.9 \\
        \quad no PG + ICP       & 4.71  & 5.09  & 142.5 & 3.3 \\
        \quad star PG + ICP     & 4.55  & 5.02  & 136.0 & 3.3 \\
        \quad mesh PG + ICP     & 5.23  & 5.69  & 135.7 & 5.2 \\
        \bottomrule
    \end{tabular}
\end{table}

\paragraph{Key findings.}
The best overall configuration for 3~robots was
\textbf{KISS-Matcher with ICP refinement and no pose graph}
($\bar{e}_t = \SI{2.04}{\metre}$, $e_t^\mathrm{RMS} =
\SI{2.37}{\metre}$).
Notably, for KISS-Matcher without ICP, increasing pose graph
connectivity improved results: mesh PG (2.85\,m~RMS) outperformed
star PG (3.19\,m) and no PG (3.44\,m).
This suggests that at the 3-robot scale, where all robots have
moderate mutual overlap, the additional constraints from the mesh
topology help average out registration noise.

FRICP consistently performed worse than KISS-Matcher across all
configurations, with rotation errors exceeding $120\degree$ on
average, indicating frequent convergence to local minima.

\paragraph{Cross-setup variability.}
Results varied substantially across position setups.
\autoref{tab:cross_setup_3r} shows the best KISS-Matcher
configuration per setup.

\begin{table}[htbp]
    \centering
    \caption{Best KISS-Matcher configuration per position setup
             (3~robots, 3 trials per configuration).}
    \label{tab:cross_setup_3r}
    \begin{tabular}{c l r r}
        \toprule
        \textbf{Setup}
            & \textbf{Best Configuration}
            & $e_t^\mathrm{RMS}$ [m]
            & $\bar{e}_r$ [deg] \\
        \midrule
        1 & mesh PG              & 0.38  &  5.0 \\
        2 & no PG                & 5.27  & 89.7 \\
        3 & star PG              & 1.74  & 53.5 \\
        4 & mesh PG + ICP        & 0.72  & 21.5 \\
        5 & no PG + ICP          & 1.48  & 13.4 \\
        \bottomrule
    \end{tabular}
\end{table}

The best-case result (\SI{0.38}{\metre}~RMS in Setup~1) was achieved
with mesh PG in a tight equilateral triangle at \SI{1.5}{\metre}
inter-robot distance (overlap: 62.3\%).
The worst case (Setup~2, \SI{5.27}{\metre}~RMS) corresponds to a
wide equilateral triangle with 55.6\% overlap, where even the best
configuration fails to achieve sub-meter accuracy.
Notably, the optimal configuration varies across setups: no single
parameter combination universally dominates.

%% ---------------------------------------------------------------------------
\section{Results: 5-Robot Experiments}
\label{sec:results_5robots}

Five homogeneous EC~Swift robots were deployed in the scalability
test world ($20 \times 20$\,m room, reduced physics) across 5 position
setups.
All 5 tests completed successfully, with test durations between
593 and \SI{656}{\second}.
The mean point cloud overlap was 48.5\% (range: 17.1--60.3\%),
notably lower than the 3-robot experiments due to larger
inter-robot distances in the wider room.

\begin{table}[htbp]
    \centering
    \caption{5-robot calibration results (averaged over 5~position
             setups, 3~trials per configuration, overlap: 48.5\%).}
    \label{tab:results_5r}
    \small
    \begin{tabular}{l r r r r}
        \toprule
        \textbf{Configuration}
            & $\bar{e}_t$ [m]
            & $e_t^\mathrm{RMS}$ [m]
            & $\bar{e}_r$ [deg]
            & $t$ [s] \\
        \midrule
        \multicolumn{5}{l}{\textit{KISS-Matcher}} \\
        \quad no PG             & \textbf{3.67}  & \textbf{4.38}  & 59.4  & 0.3 \\
        \quad star PG           & 4.16  & 5.11  & 58.0  & 0.3 \\
        \quad mesh PG           & 5.24  & 6.24  & 83.6  & 0.3 \\
        \quad no PG + ICP       & 5.49  & 6.88  & 81.9  & 2.5 \\
        \quad star PG + ICP     & 5.15  & 6.38  & 74.7  & 2.7 \\
        \quad mesh PG + ICP     & 5.16  & 6.16  & 72.6  & 5.5 \\
        \midrule
        \multicolumn{5}{l}{\textit{FRICP}} \\
        \quad no PG             & 6.24  & 6.92  & 133.3 & 4.3 \\
        \quad star PG           & 6.71  & 7.25  & 151.9 & 4.2 \\
        \quad mesh PG           & 6.85  & 7.24  & 150.0 & 11.2 \\
        \quad no PG + ICP       & 6.84  & 7.31  & 159.2 & 7.5 \\
        \quad star PG + ICP     & 6.79  & 7.32  & 160.4 & 7.6 \\
        \quad mesh PG + ICP     & 6.60  & 7.10  & 154.8 & 19.5 \\
        \bottomrule
    \end{tabular}
\end{table}

\paragraph{Key findings.}
At the 5-robot scale, \textbf{KISS-Matcher without pose graph} was
the best configuration ($e_t^\mathrm{RMS} = \SI{4.38}{\metre}$),
reversing the 3-robot trend where pose graph optimization helped.
This is because with 5~robots, the wider spatial distribution leads
to some pairwise registrations converging to incorrect local minima;
the pose graph then propagates these errors rather than correcting
them.

The 5-robot results show higher average errors than the 3-robot case
(best: 4.38\,m vs.\ 2.37\,m), primarily driven by Setup~2 (wide
pentagon, overlap 17.1\%) where all configurations produced errors
exceeding \SI{10}{\metre}.
Excluding Setup~2, the best 5-robot configuration achieves
$e_t^\mathrm{RMS} \approx \SI{2.4}{\metre}$, comparable to the
3-robot results.

The outstanding individual result was \textbf{Setup~4}
(cross/plus pattern, overlap 60.3\%), where KISS-Matcher without
pose graph achieved $e_t^\mathrm{RMS} = \SI{0.22}{\metre}$ and
$\bar{e}_r = 1.6\degree$---demonstrating that sub-meter accuracy
is achievable for 5~robots when the spatial geometry provides
sufficient feature diversity.

\paragraph{Computation time.}
KISS-Matcher completed each calibration run in
\SIrange{0.1}{0.5}{\second} at the 5-robot scale, while FRICP
required \SIrange{3}{20}{\second} (increasing significantly with
mesh topology due to $\binom{5}{2} = 10$ pairwise registrations).

%% ---------------------------------------------------------------------------
\section{Results: 10-Robot Experiments}
\label{sec:results_10robots}

Ten homogeneous EC~Swift robots were deployed in the scalability
test world across 5 position setups.
Due to limitations of the Zenoh middleware transport layer, the
10-robot experiments encountered severe reliability issues:
\begin{itemize}
    \item \textbf{Setups~3, 4, 5:} Robots failed to spawn entirely.
          The Zenoh daemon could not recover between consecutive
          tests, resulting in transport saturation errors
          (\texttt{Unable to push non droppable network message}).
    \item \textbf{Setups~1, 2:} Timed out after \SI{1000}{\second}.
          Setup~1 encountered TF frame initialization deadlocks, while
          Setup~2 completed 7 of 12 configurations before the timeout.
\end{itemize}

Due to these Zenoh transport issues, further systematic 10-robot
testing was not possible within the scope of this evaluation.
However, the partial data from Setup~2 provides limited but
informative results.

\paragraph{Partial results from Setup~2.}
The partial data from Setup~2 (wide regular decagon,
$r = \SI{5}{\metre}$, overlap: 2.1\%) includes 6 FRICP
configurations (all 3~trials each) and 1 KISS-Matcher configuration
(1~trial).
\autoref{tab:results_10r_partial} summarizes these results.

\begin{table}[htbp]
    \centering
    \caption{10-robot partial results (Setup~2 only, overlap: 2.1\%).
             Only configurations that completed at least one trial
             are shown.}
    \label{tab:results_10r_partial}
    \small
    \begin{tabular}{l r r r}
        \toprule
        \textbf{Configuration}
            & $\bar{e}_t$ [m]
            & $\bar{e}_r$ [deg]
            & $t$ [s] \\
        \midrule
        KM mesh PG + ICP   & 8.37  & 87.4  & 49.8 \\
        \midrule
        FRICP no PG        & 8.44  & 99.6  & 10.9 \\
        FRICP star PG      & 8.49  & 100.4 & 10.7 \\
        FRICP mesh PG      & 8.40  & 99.5  & 58.2 \\
        FRICP no PG + ICP  & 8.17  & 102.1 & 16.9 \\
        FRICP star PG + ICP & 8.03  & 97.2  & 17.0 \\
        \bottomrule
    \end{tabular}
\end{table}

The average translation errors of approximately \SI{8}{\metre} are
expected given the extremely low overlap (2.1\%) in the wide decagon
arrangement.
However, examining per-robot results reveals that individual robots
close to the master achieved translation errors as low as
\SI{1.6}{\milli\metre} (FRICP, trial~3), \SI{8.9}{\milli\metre}
(FRICP + ICP, trial~3), and \SI{26.8}{\milli\metre}
(FRICP star + ICP, trial~1).
These results indicate that for robot pairs with sufficient overlap,
the calibration algorithm converges correctly even in large fleets.
The high average errors are driven by distant robots with near-zero
overlap to the master, for which pairwise registration is
fundamentally infeasible.

This suggests that with a more suitable middleware configuration
(e.g., tuned Zenoh buffer sizes or alternative DDS implementations)
and tighter spatial arrangements, 10-robot calibration could achieve
substantially better results.

%% ---------------------------------------------------------------------------
\section{Realistic Proximity Scenario}
\label{sec:proximity_test}

The most common real-world use case for extrinsic multi-robot
calibration involves exactly two robots placed side by side.
A dedicated proximity test was designed to evaluate MRC under
these realistic operating conditions.

\subsection{Test Setup}

The proximity test uses the \texttt{mrc\_icp\_pg\_test\_2} launch
file, which spawns two homogeneous EC~Swift robots, each equipped
with both front and back Livox Mid-360 LiDARs (providing
$360\degree$ coverage).
The simulation uses a headless Gazebo instance with the lightweight
lidar-only world, with physics reduced to \SI{100}{\hertz} as the
robots are static.

Five predefined setups place two robots in proximity configurations
with lateral separations of \SIrange{0.8}{2.0}{\metre}:

\begin{table}[htbp]
    \centering
    \caption{Two-robot proximity test position setups.
             Robot~0 (master) is at the origin in all setups.}
    \label{tab:proximity_setups}
    \small
    \begin{tabular}{c l c c}
        \toprule
        \textbf{Setup}
            & \textbf{Description}
            & \textbf{Distance} [m]
            & $\Delta\psi$ [deg] \\
        \midrule
        1 & Parallel, no noise     & 1.5  & $0.0$ \\
        2 & Slight offset + yaw    & 1.0  & $2.9$ \\
        3 & Behind + small yaw     & 2.0  & $-4.6$ \\
        4 & Both with yaw noise    & 1.2  & $-4.6$ \\
        5 & Very close, more noise & 0.8  & $5.7$ \\
        \bottomrule
    \end{tabular}
\end{table}

\paragraph{Setup~1 anomaly.}
Setup~1 exhibited a simulation anomaly: the evaluator reported
0\% point cloud overlap despite the robots being placed at
\SI{1.5}{\metre} separation.
KISS-Matcher without ICP produced no valid registrations (0 out
of 3~trials), and all other configurations produced errors exceeding
\SI{4}{\metre}.
This anomaly is likely caused by an initialization timing issue
in the headless simulation.
Setup~1 is therefore excluded from the averaged results; the
remaining 4~setups provide sufficient coverage of the parameter
space.

In the two-robot case ($N = 2$), all three pose graph modes
(none, star, mesh) reduce to a single edge between the two robots
and are therefore expected to produce equivalent results.

\subsection{Results}

\autoref{tab:proximity_results} presents the results averaged
across Setups~2--5 (3~trials per configuration).
The mean point cloud overlap was 67.3\% (range: 64.2--70.9\%).

\begin{table}[htbp]
    \centering
    \caption{Proximity test results (2~robots, averaged over
             Setups~2--5, 3~trials per configuration,
             overlap: 67.3\%).}
    \label{tab:proximity_results}
    \small
    \begin{tabular}{l r r r}
        \toprule
        \textbf{Configuration}
            & $\bar{e}_t$ [mm]
            & $\bar{e}_r$ [deg]
            & $t$ [s] \\
        \midrule
        \multicolumn{4}{l}{\textit{KISS-Matcher}} \\
        \quad no PG             & 120.8  & 1.56  & 0.10 \\
        \quad star PG           & 153.8  & 2.08  & 0.10 \\
        \quad mesh PG           & 182.2  & 2.57  & 0.10 \\
        \quad no PG + ICP       &  76.3  & 1.02  & 0.63 \\
        \quad star PG + ICP     &  43.0  & 0.44  & 0.63 \\
        \quad mesh PG + ICP     &  26.4  & 0.24  & 0.63 \\
        \midrule
        \multicolumn{4}{l}{\textit{FRICP}} \\
        \quad no PG             &  13.0  & 0.10  & 0.78 \\
        \quad star PG           &   7.8  & 0.05  & 0.75 \\
        \quad mesh PG           & \textbf{5.4}  & \textbf{0.03}  & 0.78 \\
        \quad no PG + ICP       &  18.0  & 0.13  & 1.28 \\
        \quad star PG + ICP     &  24.3  & 0.16  & 1.30 \\
        \quad mesh PG + ICP     &  25.0  & 0.17  & 1.25 \\
        \bottomrule
    \end{tabular}
\end{table}

\paragraph{FRICP dominates in the proximity scenario.}
In stark contrast to the multi-robot results, FRICP consistently
and significantly outperformed KISS-Matcher.
The best configuration was \textbf{FRICP with mesh PG}, achieving
a mean translation error of only \SI{5.4}{\milli\metre} and a mean
rotation error of $0.03\degree$.
Even the worst FRICP configuration (mesh PG + ICP,
$\bar{e}_t = \SI{25.0}{\milli\metre}$) outperformed the best
KISS-Matcher configuration without ICP ($\bar{e}_t =
\SI{120.8}{\milli\metre}$).

This result reversal is explained by the favorable conditions
of the proximity scenario:
\begin{enumerate}
    \item \textbf{High point cloud overlap} (64--71\%) from
          $360\degree$ LiDAR coverage and short inter-robot distances.
    \item \textbf{Small relative yaw} ($< 6\degree$), so FRICP's
          initialization from the identity transform is close
          to the correct solution, avoiding local minima.
    \item \textbf{Dense geometric constraints} from both robots
          observing largely the same environment.
\end{enumerate}

\paragraph{Consistency across setups.}
\autoref{tab:proximity_per_setup} shows the best configuration's
translation error per setup, demonstrating consistent
sub-centimeter accuracy.

\begin{table}[htbp]
    \centering
    \caption{Best configuration (FRICP, mesh PG) per position setup
             (translation error averaged over 3~trials).}
    \label{tab:proximity_per_setup}
    \small
    \begin{tabular}{c c r r r}
        \toprule
        \textbf{Setup}
            & \textbf{Distance} [m]
            & $\bar{e}_t$ [mm]
            & $\bar{e}_r$ [deg]
            & \textbf{Overlap} [\%] \\
        \midrule
        2 & 1.0  &  4.5  & 0.03  & 69.1 \\
        3 & 2.0  &  5.1  & 0.02  & 64.2 \\
        4 & 1.2  &  6.4  & 0.02  & 65.0 \\
        5 & 0.8  &  5.5  & 0.03  & 70.9 \\
        \midrule
        \multicolumn{2}{c}{\textbf{Average}}
            & \textbf{5.4} & \textbf{0.03} & \textbf{67.3} \\
        \bottomrule
    \end{tabular}
\end{table}

All four valid setups achieved sub-centimeter accuracy, confirming
that MRC reliably determines the relative pose between two closely
placed robots regardless of the specific placement geometry.

%% ---------------------------------------------------------------------------
\section{Analysis: ICP Refinement Impact}
\label{sec:icp_impact}

The optional ICP refinement stage was evaluated by comparing each
backend-and-topology combination with and without ICP, both on
average across all multi-robot experiments and for the best
individual experiment.

\subsection{Average Impact (Multi-Robot)}

\autoref{tab:icp_impact} shows the average ICP impact across all
$N \in \{3, 5\}$ experiments (10~setups total).

\begin{table}[htbp]
    \centering
    \caption{ICP refinement impact on $e_t^\mathrm{RMS}$,
             averaged over all multi-robot experiments
             ($N \in \{3, 5\}$, 10~setups).
             Positive values indicate degradation.}
    \label{tab:icp_impact}
    \begin{tabular}{l r r l}
        \toprule
        \textbf{Backend + PG}
            & \textbf{Coarse} [m]
            & \textbf{+ ICP} [m]
            & \textbf{Change} \\
        \midrule
        KM, no PG    & 3.91  & 4.63  & $+18.4\%$ \\
        KM, star     & 4.15  & 4.83  & $+16.4\%$ \\
        KM, mesh     & 4.55  & 4.74  & $+4.2\%$  \\
        FRICP, no PG & 5.67  & 6.20  & $+9.3\%$  \\
        FRICP, star  & 6.00  & 6.17  & $+2.8\%$  \\
        FRICP, mesh  & 6.28  & 6.39  & $+1.8\%$  \\
        \bottomrule
    \end{tabular}
\end{table}

On average, ICP refinement degraded accuracy across all
configurations, with the largest degradation ($+18.4\%$) for
KISS-Matcher without pose graph.
The effect is smaller for configurations that already include
pose graph optimization, as the graph constraints partially
counteract the ICP-induced drift.

\subsection{Best Experiment}

Despite the negative average impact, ICP refinement provided
substantial improvements in specific setups:
\begin{itemize}
    \item \textbf{3~robots, Setup~5} (scattered asymmetric,
          overlap 51.4\%): KISS-Matcher no~PG improved from
          $e_t^\mathrm{RMS} = \SI{4.27}{\metre}$ to
          $\SI{1.48}{\metre}$ ($-65.3\%$) with ICP.
    \item \textbf{5~robots, Setup~1} (regular pentagon,
          overlap 55.3\%): KISS-Matcher mesh~PG improved from
          $e_t^\mathrm{RMS} = \SI{3.13}{\metre}$ to
          $\SI{1.98}{\metre}$ ($-36.7\%$) with ICP.
\end{itemize}

These cases share a common characteristic: the coarse registration
produced a moderately correct alignment (within a few meters),
providing ICP with a sufficiently good initial guess to converge
toward the global minimum.
In setups where the coarse registration is already very good
($< 1$\,m) or very bad ($> 5$\,m), ICP either has nothing to
improve or diverges to a local minimum.

\subsection{Proximity Scenario}

In the proximity scenario, ICP refinement exhibited the opposite
behavior for the two backends:
\begin{itemize}
    \item \textbf{KISS-Matcher:} ICP improved accuracy significantly,
          reducing translation error by 37\% (no~PG), 72\% (star),
          and 86\% (mesh).
          The high overlap and dense point clouds ensure sufficient
          correspondences for ICP convergence.
    \item \textbf{FRICP:} ICP \emph{degraded} accuracy by
          38--363\%, as FRICP already converges to a near-perfect
          solution; the subsequent ICP pass introduces systematic
          bias from the non-uniform LiDAR point distribution.
\end{itemize}

%% ---------------------------------------------------------------------------
\section{Analysis: Pose Graph Optimization}
\label{sec:pose_graph_impact}

The impact of pose graph optimization was evaluated by comparing
three modes: no pose graph (direct pairwise registration), star
topology ($N - 1$ edges), and mesh topology ($\binom{N}{2}$ edges).

\subsection{Average Impact (Multi-Robot)}

\autoref{tab:pg_impact} shows the average pose graph impact across
all multi-robot experiments ($N \in \{3, 5\}$, 10~setups).

\begin{table}[htbp]
    \centering
    \caption{Pose graph optimization impact on $e_t^\mathrm{RMS}$,
             averaged over all multi-robot experiments.
             Positive values indicate degradation vs.\ no~PG.}
    \label{tab:pg_impact}
    \begin{tabular}{l r r r}
        \toprule
        \textbf{Backend}
            & \textbf{No PG} [m]
            & \textbf{Star} [m]
            & \textbf{Mesh} [m] \\
        \midrule
        KM           & 3.91  & 4.15 ($+6.1\%$)  & 4.55 ($+16.4\%$)  \\
        KM + ICP     & 4.63  & 4.83 ($+4.3\%$)  & 4.74 ($+2.4\%$)   \\
        FRICP        & 5.67  & 6.00 ($+5.8\%$)  & 6.28 ($+10.8\%$)  \\
        FRICP + ICP  & 6.20  & 6.17 ($-0.5\%$)  & 6.39 ($+3.1\%$)   \\
        \bottomrule
    \end{tabular}
\end{table}

When averaged across all setups and fleet sizes, pose graph
optimization generally degraded calibration accuracy.
For KISS-Matcher, star topology increased the RMS error by 6.1\%
and mesh topology by 16.4\%.
This occurs because in setups where some pairwise registrations
converge to incorrect local minima, the pose graph optimizer
propagates these errors through the graph rather than correcting
them.

The only near-neutral case was FRICP with star~PG and ICP
($-0.5\%$), where the combination of dense point-to-point
optimization and the pose graph noise model marginally improved
robustness.

\subsection{Best Experiment}

In specific setups, pose graph optimization provided substantial
improvements:
\begin{itemize}
    \item \textbf{3~robots, Setup~4} (structured composite,
          overlap 59.3\%): KISS-Matcher mesh~PG improved from
          $e_t^\mathrm{RMS} = \SI{4.84}{\metre}$ (no~PG) to
          $\SI{0.95}{\metre}$ ($-80.4\%$).
    \item \textbf{3~robots, Setup~1} (tight triangle, overlap 62.3\%):
          KISS-Matcher star~PG improved from
          $e_t^\mathrm{RMS} = \SI{0.77}{\metre}$ (no~PG) to
          $\SI{0.39}{\metre}$ ($-49.4\%$).
\end{itemize}

These best cases occur when all pairwise registrations in the
graph converge to correct solutions; the graph optimizer then
reduces residual noise.
In setups with even one badly converged edge, however, the graph
propagates the error, explaining the negative average effect.

\subsection{Proximity Scenario}

In the proximity scenario, pose graph topology had a differential
impact:
\begin{itemize}
    \item \textbf{FRICP:} The pose graph improved accuracy
          from \SI{13.0}{\milli\metre} (no~PG) to
          \SI{7.8}{\milli\metre} (star, $-40\%$) and
          \SI{5.4}{\milli\metre} (mesh, $-58\%$).
    \item \textbf{KISS-Matcher without ICP:} The pose graph
          degraded results from \SI{120.8}{\milli\metre} (no~PG)
          to \SI{182.2}{\milli\metre} (mesh, $+51\%$).
    \item \textbf{KISS-Matcher with ICP:} The pose graph improved
          results from \SI{76.3}{\milli\metre} (no~PG + ICP) to
          \SI{26.4}{\milli\metre} (mesh + ICP, $-65\%$).
\end{itemize}

The fact that pose graph consistently helps FRICP in the
proximity scenario (where all pairwise registrations are correct)
but hurts in the multi-robot scenario (where some edges are
unreliable) confirms that the benefit of graph optimization
depends fundamentally on edge quality.

%% ---------------------------------------------------------------------------
\section{Summary}
\label{sec:evaluation_summary}

The experimental evaluation comprised over 700 individual
calibration runs across 2 deployment scenarios, 3 fleet sizes,
15 position setups, and 12 parameter combinations.
The key findings are:

\begin{enumerate}
    \item \textbf{The optimal backend depends on the deployment
          scenario.}
          In the multi-robot distributed scenario ($N \geq 3$),
          KISS-Matcher consistently outperforms FRICP, which
          frequently converges to local minima (rotation errors
          $> 120\degree$).
          In the two-robot proximity scenario, the result reverses:
          FRICP achieves sub-centimeter accuracy
          ($\bar{e}_t = \SI{5.4}{\milli\metre}$) while KISS-Matcher
          is limited to $\sim$\,\SI{12}{\centi\metre} by its
          feature quantization.

    \item \textbf{The optimal parameter combination is
          fleet-size-dependent.}
          For 3~robots, KISS-Matcher with ICP and no pose graph is
          best ($e_t^\mathrm{RMS} = \SI{2.37}{\metre}$).
          For 5~robots, KISS-Matcher without ICP or pose graph is
          best ($e_t^\mathrm{RMS} = \SI{4.38}{\metre}$).
          For 2~robots in proximity, FRICP with mesh PG is best
          ($\bar{e}_t = \SI{5.4}{\milli\metre}$).

    \item \textbf{ICP refinement and pose graph optimization are
          scenario-dependent.}
          Averaged across all multi-robot experiments, both ICP
          refinement ($+2$--$18\%$ RMS) and pose graph optimization
          ($+2$--$16\%$ RMS) degraded accuracy.
          However, in favorable setups (high overlap, correct coarse
          alignment), ICP improved accuracy by up to 65\% and mesh
          PG by up to 80\%.
          In the proximity scenario, ICP dramatically improved
          KISS-Matcher ($-86\%$) but degraded FRICP ($+363\%$).

    \item \textbf{Sub-centimeter accuracy is achievable for the
          primary use case.}
          The two-robot proximity scenario achieved
          $\bar{e}_t = \SI{5.4}{\milli\metre}$ and
          $\bar{e}_r = 0.03\degree$ with FRICP and mesh PG,
          consistent across all four valid position setups
          (\SIrange{0.8}{2.0}{\metre} separation).

    \item \textbf{Results are sensitive to geometric context.}
          Calibration accuracy varies significantly across spatial
          arrangements.
          The best 5-robot result
          ($e_t^\mathrm{RMS} = \SI{0.22}{\metre}$, Setup~4)
          demonstrates that sub-meter accuracy is achievable when
          the geometry provides diverse LiDAR features.

    \item \textbf{Scalability beyond 5~robots requires middleware
          improvements.}
          The 10-robot tests were limited by Zenoh middleware
          transport saturation, preventing systematic evaluation.
          The partial results show that nearby robots achieve
          sub-centimeter calibration even in large fleets,
          suggesting that the MRC algorithm itself scales, but the
          communication infrastructure requires tuning for
          large fleet sizes.

    \item \textbf{Recommended configurations:}
          \begin{itemize}
              \item \textbf{Two-robot proximity:} FRICP with mesh PG
                    ($\bar{e}_t \approx \SI{5}{\milli\metre}$,
                    $t \approx \SI{0.8}{\second}$).
              \item \textbf{3-robot distributed:} KISS-Matcher with
                    ICP, no PG
                    ($e_t^\mathrm{RMS} \approx \SI{2.4}{\metre}$,
                    $t \approx \SI{1.2}{\second}$).
              \item \textbf{5+ robot distributed:} KISS-Matcher
                    without ICP or PG
                    ($e_t^\mathrm{RMS} \approx \SI{4.4}{\metre}$,
                    $t \approx \SI{0.3}{\second}$).
          \end{itemize}
\end{enumerate}
